\documentclass[10pt]{beamer}
\usetheme{Madrid}
\usecolortheme{whale}
\usepackage[utf8]{inputenc}
\usepackage[T1]{fontenc}
\usepackage[spanish,es-tabla,es-nodecimaldot]{babel}
\usepackage{amsmath,amssymb}
\usepackage{booktabs}
\usepackage{array}
\newcolumntype{L}[1]{>{\raggedright\arraybackslash}p{#1}}
\usepackage{tikz}
\usetikzlibrary{babel}

\setbeamertemplate{navigation symbols}{}
\setbeamerfont{caption}{size=\scriptsize}
\newcommand{\dato}[1]{\textbf{\color{blue!60!black}#1}}

\title[Laguna Torca: resultados]{Superficie hídrica de la Laguna Torca}
\subtitle{Resultados de los modelos de predicción}
\author[J. P. Maureira]{Juan Pablo Maureira}
\institute[UTalca]{Magíster en Data Science\\Universidad de Talca}
\date{\today}

\begin{document}

\begin{frame}
\titlepage
\end{frame}

% ------------------------------------------------------------
\begin{frame}{Datos usados}
\begin{itemize}
\item \textbf{Superficie:} 959 escenas Landsat 5, 7, 8 y 9 (1985--2024), índice NDWI en Google Earth Engine, promedio mensual sin imputar.
\item \textbf{Clima histórico:} CR2MET v2.0 en el píxel de la laguna (precipitación, temperatura máxima y mínima, 1979--2019).
\item \textbf{Proyecciones:} Explorador Climático CR2, RCP~2.6 (26 modelos globales) y RCP~8.5 (39 modelos), hasta 2095.
\item \textbf{Validación:} siempre temporal (se entrena con el pasado y se prueba con el futuro) y contra referencias simples: media, climatología y persistencia.
\end{itemize}
\end{frame}

% ------------------------------------------------------------
\begin{frame}{La serie observada: poca variación y un salto de sensor}
\begin{columns}[T]
\begin{column}{0.56\textwidth}
\centering
\begin{tikzpicture}[x=0.68cm,y=0.95cm,font=\scriptsize]
\draw[gray!50] (0,0) rectangle (8,3.5);
\foreach \v/\l in {0/110,1/120,2/130,3/140} {\draw[gray!30] (0,\v)--(8,\v); \node[left] at (0,\v) {\l};}
\foreach \x/\l in {0/1985,2/1995,4/2005,6/2015,8/2025} {\node[below] at (\x,0) {\l};}
\draw[red!70!black,dashed] (5.4,0)--(5.4,3.5);
\node[red!70!black,above] at (5.4,3.5) {cambio de sensor};
\draw[blue!70!black,thick] plot coordinates {(0.20,1.923) (0.40,1.477) (0.60,2.082) (0.80,1.936) (1.00,1.598) (1.20,1.716) (1.80,0.880) (2.00,1.564) (2.20,1.247) (2.40,1.717) (2.60,0.615) (2.80,1.065) (3.00,0.825) (3.20,1.168) (3.40,1.336) (3.60,1.191) (3.80,1.425) (4.00,1.660) (4.20,1.692) (4.40,1.555) (4.60,0.689) (4.80,1.429) (5.00,1.144) (5.20,1.561)};
\draw[orange!80!black,thick] plot coordinates {(5.60,2.320) (5.80,2.357) (6.00,1.776) (6.20,2.297) (6.40,2.453) (6.60,2.694) (6.80,2.015) (7.00,2.549) (7.20,3.155) (7.40,2.675) (7.60,3.168) (7.80,3.078)};
\node[rotate=90] at (-1.3,1.75) {Área media anual (ha)};
\node[blue!70!black] at (2.6,3.0) {Landsat 5/7};
\node[orange!80!black] at (6.7,1.2) {Landsat 8/9};
\end{tikzpicture}
\end{column}
\begin{column}{0.42\textwidth}
\small
\begin{itemize}
\item Área anual entre \dato{116 y 142 ha}.
\item Media Landsat 5/7: \dato{124.4 ha}; Landsat 8/9: \dato{135.5 ha}.
\item Salto de \dato{11.0 ha} en 2012 (8.7 ha descontando la lluvia), justo en el cambio de sensor.
\item En pares de escenas simultáneas, Landsat 7 mide \dato{3.2 ha menos} que Landsat 5.
\end{itemize}
\end{column}
\end{columns}
\vspace{2mm}
\small Lo que hay que explicar es una variación de pocas hectáreas, del mismo orden que el error de medición.
\end{frame}

% ------------------------------------------------------------
\begin{frame}{Correlación con el clima}
\centering
\small
\begin{tabular}{llrr}
\toprule
Escala & Variable con mayor asociación & $\rho$ de Spearman & $R^2$ aprox. \\
\midrule
Mensual (208 meses) & Temperatura máxima & $-0.21$ & 4.3\,\% \\
Mensual (208 meses) & Precipitación acumulada 12 meses & $0.16$ & 2.4\,\% \\
Anual (31 años) & Precipitación acumulada 36 meses & $0.27$ ($p=0.14$) & 7.3\,\% \\
Anual (31 años) & SPEI-12 & $0.21$ ($p=0.25$) & 4.6\,\% \\
\bottomrule
\end{tabular}
\vspace{3mm}
\begin{itemize}
\item Correlaciones sin ciclo estacional y sin el salto de sensor.
\item La correlación ``cruda'' llegaba a 0.32, pero venía del ciclo estacional y del cambio de sensor, no del clima.
\item Ni con ventanas largas de lluvia (12 a 36 meses) ni con SPEI aparece una relación significativa a escala anual.
\end{itemize}
\end{frame}

% ------------------------------------------------------------
\begin{frame}{Modelos mensuales: ninguno supera a la persistencia}
\centering
\small
\begin{tabular}{lrr}
\toprule
Modelo (validación temporal, 4 pliegues) & RMSE (ha) & MAE (ha) \\
\midrule
SARIMAX(1,0,0) sin clima & 7.68 & 5.32 \\
Ridge sin clima (estacionalidad + época) & 7.70 & 5.60 \\
Persistencia: último valor (referencia) & 7.82 & 5.65 \\
Random Forest con clima & 7.86 & 5.69 \\
Ridge con clima & 7.91 & 5.74 \\
XGBoost con clima & 8.37 & 6.09 \\
SARIMAX(1,0,0) con clima & 14.07 & 12.22 \\
Media histórica (referencia) & 9.38 & 7.27 \\
\bottomrule
\end{tabular}
\vspace{3mm}
\begin{itemize}
\item Diferencias de 0.1 a 0.2 ha, frente a una desviación entre pliegues de ${\approx}2$ ha: es un \textbf{empate}.
\item Agregar clima no mejora ningún modelo (Ridge $+0.21$ ha, XGBoost $+0.33$ ha, SARIMAX $+6.4$ ha).
\item Lo que gana sobre la media es la variable de \textbf{época de sensor}, no el clima.
\end{itemize}
\end{frame}

% ------------------------------------------------------------
\begin{frame}{El modelo anual con imputación MICE}
\textbf{Es el único modelo que separa claramente RCP~2.6 de RCP~8.5.}
\vspace{2mm}
\begin{enumerate}
\item \textbf{Imputación:} MICE con Random Forest rellena los 212 meses sin escena usando el clima del mismo mes.
\item \textbf{Suavizado:} filtro Savitzky--Golay (13 meses, grado 3) sobre la serie completa.
\item \textbf{Serie objetivo:} promedio anual de esa serie, 1987--2019.
\item \textbf{Predictores:} precipitación, temperatura máxima y mínima anuales, con rezagos (7 variables).
\item \textbf{Modelo:} regresión lineal, proyectada con cada modelo global (media móvil de 3 años).
\end{enumerate}
\vspace{2mm}
Se evaluaron tres versiones: \textbf{original} (como en el notebook inicial), \textbf{corregida} (lectura correcta de los archivos CR2) y \textbf{corregida + época de sensor}.
\end{frame}

% ------------------------------------------------------------
\begin{frame}{Modelo MICE: la diferencia entre escenarios crece con el horizonte}
\centering
\small
\begin{tabular}{llrrr}
\toprule
Versión & Horizonte & RCP~2.6 (ha) & RCP~8.5 (ha) & RCP~8.5 $-$ RCP~2.6 \\
\midrule
Original & 2025--2040 & 97.8 & 101.3 & \dato{3.6} \\
Original & 2041--2070 & 99.0 & 109.8 & \dato{10.8} \\
Original & 2071--2095 & 99.6 & 122.2 & \dato{22.6} \\
\midrule
Corregida & 2025--2040 & 131.7 & 134.0 & 2.2 \\
Corregida & 2041--2070 & 134.0 & 142.3 & 8.3 \\
Corregida & 2071--2095 & 134.4 & 153.9 & 19.5 \\
\midrule
Corregida + época & 2025--2040 & 133.1 & 133.1 & 0.1 \\
Corregida + época & 2041--2070 & 133.1 & 133.1 & 0.0 \\
Corregida + época & 2071--2095 & 133.4 & 133.4 & 0.0 \\
\bottomrule
\end{tabular}
\vspace{3mm}
\begin{itemize}
\item Bajo RCP~8.5 (más seco y más cálido) el modelo predice \textbf{más} agua.
\item Al agregar la época de sensor, la diferencia \textbf{desaparece}.
\end{itemize}
\end{frame}

% ------------------------------------------------------------
\begin{frame}{MICE comparado con los demás modelos}
\begin{columns}[T]
\begin{column}{0.47\textwidth}
\centering
\scriptsize RCP~8.5 $-$ RCP~2.6 en 2071--2095 (ha)\\[1mm]
\begin{tikzpicture}[x=1cm,y=1cm,font=\scriptsize]
\draw[gray!60] (0,-0.3)--(0,4.3);
\foreach \v in {0,10,20} {\draw[gray!30] (\v*0.13,-0.3)--(\v*0.13,4.3); \node[below] at (\v*0.13,-0.3) {\v};}
\fill[red!65!black] (0,3.8) rectangle (22.6*0.13,4.2); \node[left] at (0,4.0) {MICE original};
\node[right] at (22.6*0.13,4.0) {22.6};
\fill[red!40] (0,3.0) rectangle (19.5*0.13,3.4); \node[left] at (0,3.2) {MICE corregido};
\node[right] at (19.5*0.13,3.2) {19.5};
\fill[gray] (0,2.2) rectangle (0.04,2.6); \node[left] at (0,2.4) {MICE + época};
\node[right] at (0.05,2.4) {0.0};
\fill[blue!60!black] (0,1.4) rectangle (0.14*0.13,1.8); \node[left] at (0,1.6) {Anual observado};
\node[right] at (0.05,1.6) {0.1};
\fill[blue!60!black] (-0.6*0.13,0.6) rectangle (0,1.0); \node[left] at (-0.1,0.8) {Mensual observado};
\node[right] at (0.02,0.8) {$-0.6$};
\end{tikzpicture}
\end{column}
\begin{column}{0.5\textwidth}
\scriptsize
\begin{tabular}{lr}
\toprule
Validación (área anual observada) & RMSE \\
\midrule
Año anterior (referencia) & \textbf{4.51} \\
MICE + época & 4.80 \\
Media histórica (referencia) & 5.91 \\
MICE original & 6.02 \\
\bottomrule
\end{tabular}
\vspace{2mm}

\small
\begin{itemize}
\item El modelo MICE original \textbf{pierde contra la media} y contra repetir el año anterior.
\item Los modelos con área observada predicen una laguna estable (${\approx}129$--$133$ ha) en ambos escenarios.
\end{itemize}
\end{column}
\end{columns}
\end{frame}

% ------------------------------------------------------------
\begin{frame}{¿Por qué el modelo MICE separa los escenarios?}
\begin{enumerate}
\item \textbf{Circularidad de la imputación.} En los meses imputados el área correlaciona con la lluvia ($\rho=0.27$, $p<0.001$); en los observados, no ($\rho=0.06$). MICE le ``inyecta'' clima a la serie.
\item \textbf{Temperaturas que se compensan.} En la versión corregida, de los 19.5 ha de diferencia, la precipitación aporta \dato{2.8 ha}; la mínima suma \dato{$+35.4$ ha} y la máxima resta \dato{18.7 ha}. Son variables correlacionadas con coeficientes de signo opuesto.
\item \textbf{Extrapolación.} Bajo RCP~8.5 las temperaturas quedan unos 2\,$^\circ$C por sobre todo lo visto en el entrenamiento.
\item \textbf{Confusión con el sensor.} El área sube en 2013 por el cambio de Landsat, y en esos años también sube la temperatura. Con la variable de época, la diferencia cae a \dato{0.0 ha}.
\end{enumerate}
\vspace{2mm}
\centering
\fbox{\parbox{0.9\textwidth}{\centering La separación entre escenarios es un artefacto del método, no una señal climática.}}
\end{frame}

% ------------------------------------------------------------
\begin{frame}{Comparación con la literatura}
\centering
\scriptsize
\setlength{\tabcolsep}{4pt}
\begin{tabular}{L{3.3cm}L{1.9cm}L{3.1cm}L{2.5cm}}
\toprule
Estudio & Cambio de área & Resultado & Validación \\
\midrule
Wang et al. (2023), 8 modelos de ML & n.d. & RF + SPEI-9: $R^2=0.88$ & División aleatoria 90/10 \\
Lago Chungará & n.d. & Lineal lluvia--área: $R^2\approx0.80$ & Ajuste descriptivo \\
Bakhtegan y Tashk, Irán (2023) & Se secaron & Lluvia $r\approx0.55$; ETP $r\approx-0.68$ & Proyección RCP~4.5/8.5 \\
Fuentealba et al. (2021), Chile & $-7$ a $-25$\,\% & Responden a la megasequía & Correlación descriptiva \\
\midrule
\textbf{Torca, modelo MICE} & ${\approx}\pm10$\,\% & Separa escenarios (22.6 ha) & Pierde contra la media \\
\textbf{Torca, área observada} & ${\approx}\pm10$\,\% & $|\rho|\le0.27$; $R^2\le7$\,\% & Temporal + persistencia \\
\bottomrule
\end{tabular}
\vspace{3mm}
\small
\begin{itemize}
\item Los casos con relación clara (Bakhtegan, Aculeo, lagos andinos) tienen \textbf{grandes cambios de superficie} y usan ventanas de lluvia largas (SPEI-9 a 12, acumulados de años).
\item Varios usan división aleatoria, que infla el $R^2$ en series autocorrelacionadas.
\item Cerca del 11\,\% del agua continental cae en píxeles mixtos de 30 m (Pickens et al., 2020): en una laguna de 130 ha, solo el borde mueve la medición varias hectáreas.
\end{itemize}
\end{frame}

% ------------------------------------------------------------
\begin{frame}{Conclusiones}
\begin{itemize}
\item La Torca varía muy poco (116--142 ha), y el único quiebre coincide con el cambio de sensor.
\item El clima explica a lo más un 4--7\,\% de la variación, y ningún modelo con clima supera a la persistencia.
\item El modelo MICE es el único que diferencia RCP~2.6 de RCP~8.5 (hasta \textbf{22.6 ha}), pero:
\begin{itemize}
\item pierde contra la media en validación,
\item predice más agua en el escenario más seco,
\item y la diferencia desaparece al controlar por el sensor.
\end{itemize}
\item Con los modelos validados, la proyección a 2095 es una laguna \textbf{estable en ${\approx}129$--$133$ ha} en ambos escenarios, con un error de ${\approx}5$--$8$ ha.
\item A diferencia de los casos de la literatura, el área de la Torca no parece regida por el clima local: posiblemente la regulan su conexión con el sistema lacustre, el manejo o el agua subterránea.
\end{itemize}
\end{frame}

% ------------------------------------------------------------
\begin{frame}{Referencias}
\scriptsize
\begin{itemize}
\item Fuentealba et al. (2021). Respuesta de lagos de Chile central a la megasequía.\\ \url{https://doaj.org/article/432f44b4ace347a88a118a966a7fd5e0}
\item Pickens et al. (2020). Mapping and sampling to characterize global inland water dynamics from 1999 to 2018 with full Landsat time-series. \emph{Remote Sensing of Environment}.
\item Shrestha et al. (2021). Lagos de los Sand Hills de Nebraska y sequía.\\ \url{https://doaj.org/article/d381c132fca242daa33c03819a50009c}
\item Wang et al. (2023). Comparison of machine learning models to predict lake area. \emph{Remote Sensing}, 15, 4153.
\item Lagos Bakhtegan y Tashk (2023). \emph{The Egyptian Journal of Remote Sensing and Space Sciences}.\\ \url{https://doaj.org/article/5b53145a0bd84e99a6596e6f6635e808}
\item Lago Chungará: relación lluvia--área. \emph{Caminhos de Geografia}.\\ \url{https://seer.ufu.br/index.php/caminhosdegeografia/article/download/75798/41275/365954}
\end{itemize}
\end{frame}

\end{document}
